\documentclass[a4paper,12pt]{article}

\usepackage[T1]{fontenc}
\usepackage[utf8]{inputenc}
\usepackage[ngerman]{babel}
\usepackage{csquotes}
\usepackage{enumitem}
\usepackage{parskip}

\title{Visible Light Communication\\[0.4em]
  \large Tiefpasseffekte und Konsequenzen für die PPM-Modulation}
\author{}
\date{}

\begin{document}

\maketitle

Dein beobachtetes Verhalten ist absolut typisch für Visible Light Communication (VLC)
und lässt sich sehr gut aus Sicht der Hochfrequenz- und Systemtechnik erklären.

% ------------------------------------------------------------------
\section*{Warum dein Rechteck zu \enquote{Haifischflossen} wird}

In deiner gesamten Übertragungsstrecke wirken mehrere Komponenten wie \textbf{Tiefpassfilter}:

\begin{itemize}
  \item Treiberschaltung (Transistor + hohe Spannung)
  \item LED/Leuchte
  \item Photodiode + TIA (Transimpedanzverstärker)
  \item Leitungen und parasitäre Kapazitäten
\end{itemize}

Ein ideales Rechtecksignal besteht aus \textbf{sehr vielen hohen Frequenzanteilen} (Oberwellen).
Ein Tiefpass dämpft genau diese hohen Frequenzen.

\textbf{Ergebnis:}
\begin{itemize}
  \item Flanken werden \textbf{abgerundet}
  \item Anstiege werden \textbf{langsamer}
  \item Statt Rechteck bekommst du exponentielle Anstiege/Abfälle
        $\rightarrow$ \enquote{Haifischflossen}
\end{itemize}

% ------------------------------------------------------------------
\section*{Was bedeutet das konkret für PPM (Pulse Position Modulation)?}

Bei \textbf{PPM} ist \textbf{nicht die Amplitude entscheidend}, sondern:

\begin{quote}
  $\rightarrow$ \textbf{der exakte Zeitpunkt eines Pulses}
\end{quote}

Und genau hier entsteht dein Problem.

% ------------------------------------------------------------------
\section*{Konsequenzen der Tiefpass-Effekte für PPM}

\subsection*{1.\ Zeitliche Verschiebung (Jitter\,/\,Delay)}

Durch die langsamen Flanken:

\begin{itemize}
  \item Der Puls erreicht den Schwellwert \textbf{später}
  \item $\rightarrow$ Der gemessene Zeitpunkt ist \textbf{verschoben}
\end{itemize}

Bei PPM führt das direkt zu:

\begin{itemize}
  \item falscher Symbolinterpretation
  \item Bitfehlern
\end{itemize}

% ------------------------------------------------------------------
\subsection*{2.\ Pulsverbreiterung (Inter-Symbol Interference, ISI)}

Da dein Signal nicht sofort wieder auf~0 fällt:

\begin{itemize}
  \item Ein Puls \enquote{zieht sich} in das nächste Zeitfenster hinein
  \item $\rightarrow$ Symbole beeinflussen sich gegenseitig
\end{itemize}

Besonders kritisch bei:

\begin{itemize}
  \item hohen Datenraten
  \item kurzen Symbolzeiten
\end{itemize}

% ------------------------------------------------------------------
\subsection*{3.\ Amplitudenreduktion}

Durch die Filterung:

\begin{itemize}
  \item Pulse werden kleiner
  \item $\rightarrow$ schlechteres Signal-Rausch-Verhältnis (SNR)
\end{itemize}

Folge:

\begin{itemize}
  \item Schwellwert-Detektion wird unsicher
\end{itemize}

% ------------------------------------------------------------------
\subsection*{4.\ Verlust der Pulsform}

PPM geht davon aus:

\begin{itemize}
  \item klare, schmale Pulse
\end{itemize}

Tatsächlich bekommst du:

\begin{itemize}
  \item asymmetrische, verzerrte Pulse (\enquote{Haifischflossen})
\end{itemize}

Das erschwert:

\begin{itemize}
  \item Synchronisation
  \item saubere Slot-Erkennung
\end{itemize}

% ------------------------------------------------------------------
\section*{Intuitive Zusammenfassung}

Dein System verhält sich wie:

\begin{quote}
  Ein langsames System versucht, schnelle Information zu übertragen.
\end{quote}

Und PPM ist besonders empfindlich, weil:

\begin{itemize}
  \item \textbf{Zeitinformation} übertragen wird
  \item nicht nur \enquote{Signal da\,/\,nicht da}
\end{itemize}

% ------------------------------------------------------------------
\section*{Was heißt das praktisch für dich?}

\subsection*{Ohne Gegenmaßnahmen}

\begin{itemize}
  \item steigende Bitfehlerrate (BER)
  \item limitierte maximale Datenrate
  \item instabile Dekodierung
\end{itemize}

% ------------------------------------------------------------------
\section*{Typische Lösungen\,/\,Gegenmaßnahmen}

\subsection*{1.\ Bandbreite erhöhen}

\begin{itemize}
  \item schnellere LED\,/\,Treiber
  \item schnellerer TIA
  \item kleinere parasitäre Kapazitäten
\end{itemize}

\subsection*{2.\ Entzerrung (Equalization)}

\begin{itemize}
  \item Vorverzerrung am Sender (Pre-Emphasis)
  \item oder digitale Nachbearbeitung
\end{itemize}

\subsection*{3.\ Angepasste Modulation}

\begin{itemize}
  \item größere Zeitabstände bei PPM
  \item oder robustere Verfahren:
    \begin{itemize}
      \item OOK (On-Off Keying)
      \item DCO-OFDM (in komplexeren VLC-Systemen)
    \end{itemize}
\end{itemize}

\subsection*{4.\ Schwellwert-Optimierung}

\begin{itemize}
  \item adaptive Thresholds statt fixer Werte
\end{itemize}

% ------------------------------------------------------------------
\section*{Fazit}

Die \enquote{Haifischflossen} sind kein Fehler, sondern ein Hinweis:

\begin{quote}
  $\rightarrow$ \textbf{Dein System ist bandbreitenbegrenzt.}
\end{quote}

Für PPM bedeutet das konkret:

Die Information (Zeitposition der Pulse) wird durch Verzögerung, Verbreiterung und Verzerrung
verfälscht -- was direkt zu Dekodierfehlern führt.

\end{document}
